% Drawing 142/06-8 -- folder 142, batch 06, archive page 109 (PDF page 110).
% Transcription: Partie retournée. Une sphère avec pôles $0$ et $\infty$, l'équateur passant par $1$, quatre petits cercles sur l'équateur ($a$, $b$ marqués) ; deux flèches en partent vers le bas, avec la légende « division par $D$ (multiplication par les \ill{} \uncertain{$R_s$}) ». Une seconde flèche « division par $D^+$ (\uncertain{multiplication} \ill{} \uncertain{sur} $(0, \infty)$) » mène à une sphère dont l'équateur passe par $1$.
% Kind: spheres with marked circles and arrows (upside-down half of the sheet). Method 2 (schematic redraw from the facsimile). Status: partial.
% Notes: Upside-down half of p. 109, redrawn upright. Left: sphere with poles 0 (top) and infinity (bottom), dashed central meridian, the equator through 1, four small circles near the equator (a, b marked); two arrows go down from it, with the legend "division par D (multiplication par les ... R_s)" (drawn as text, the illegible word left out). Right: the arrow "division par D^+ (...)" and the sphere at the top right, with the pole circle infinity and the inked equator passing through 1 (ringed); on the page the D^+ arrow points rather to the central sphere (06-9) than to this one. Small labels and the faint pencil meridians are approximate.
% Coordinates: 1 unit = 1 mm = 1/100 of the crop width (bbox in page fractions: [0.196, 0.563, 0.884, 0.88]).
% Redrawn by Claude (claude-opus-5-5), issue #73, 2026-11-07; not checked by a human; not compiled (no TeX in the session).
% Rights: reproduces part of the work -- behind RENDER_TRANSCRIPTS, never published by itself.
\documentclass[tikz,border=3pt]{standalone}
\usepackage{amsmath,amssymb}
\usetikzlibrary{patterns,arrows.meta}
\tikzset{pencil/.style={black!45}, struck/.style={black!35}, mark/.style={circle,draw,inner sep=1.3pt}, pt/.style={circle,fill,inner sep=0.9pt}, opt/.style={circle,draw,fill=white,inner sep=1.1pt}}
\begin{document}
\begin{tikzpicture}[x=1mm,y=1mm,line cap=round,line join=round,font=\scriptsize]
\draw[pencil] (23.3,-24.4) ellipse[x radius=14, y radius=15.1];
\node[pt] at (23.3,-7) {};
\node[pt] at (24.4,-40.7) {};
\draw (23.3,-10.5) ellipse[x radius=2.3, y radius=0.9];
\draw (24.4,-38.4) ellipse[x radius=2.3, y radius=0.9];
\draw[dashed] (23.3,-11.6) -- (24.4,-38.4);
\draw[pencil] plot[smooth, tension=0.55] coordinates {(20.9,-10.5) (12.8,-17.4) (10.5,-25.6) (15.1,-36) (22.1,-38.4)};
\draw[pencil] plot[smooth, tension=0.55] coordinates {(25.6,-10.5) (33.7,-17.4) (36,-25.6) (31.4,-36) (26.7,-38.4)};
\draw (23.3,-25.6) ellipse[x radius=14, y radius=2.6];
\draw (17.4,-22.1) circle[radius=1.4];
\draw (30.2,-20.9) circle[radius=1.4];
\draw (17.4,-26.7) circle[radius=1.4];
\draw (26.7,-26.2) circle[radius=1.7];
\node[inner sep=1pt] at (39.5,-23.3) {\scriptsize $1$};
\node[inner sep=1pt] at (26.7,-12.8) {\tiny $a$};
\node[inner sep=1pt] at (19.8,-28.5) {\tiny $b$};
\node[inner sep=1pt] at (23.3,-3.5) {\tiny $0$};
\node[inner sep=1pt] at (24.4,-44.2) {\tiny $\infty$};
\draw[-{Stealth[length=2mm]}] (10.5,-43) -- (2.3,-60.5);
\draw[-{Stealth[length=2mm]}] (17.4,-40.7) -- (12.8,-55.8);
\node[inner sep=1pt] at (23.3,-51.2) {\tiny division par $D$};
\node[inner sep=1pt] at (23.3,-55.8) {\tiny (multiplication};
\node[inner sep=1pt] at (23.3,-59.3) {\tiny par les \dots\ $R_s$)};
\draw[-{Stealth[length=2mm]}] (41.9,-29.1) -- (54.7,-54.7);
\node[inner sep=1pt] at (51.2,-41.9) {\tiny division par $D^+$};
\node[inner sep=1pt] at (51.2,-46.5) {\tiny (multiplication \dots};
\node[inner sep=1pt] at (84.9,-40.7) {\tiny sur $(0,\infty)$)};
\draw (76.7,-30.2) circle[radius=18.6];
\draw (77.9,-15.1) ellipse[x radius=2.9, y radius=1.2];
\node[inner sep=1pt] at (79.1,-11.6) {\tiny $\infty$};
\draw[very thick] (75.6,-30.2) ellipse[x radius=15.1, y radius=4.7];
\node[mark] at (91.9,-31.4) {};
\node[inner sep=1pt] at (95.3,-31.4) {\scriptsize $1$};
\end{tikzpicture}
\end{document}
